% Appendix A - Validaciones Metodológicas

\chapter{Validaciones Metodológicas}
\label{ch:appendix_validaciones}

\section{Validación de la Implementación de SMOTE}
\label{sec:validacion_smote}

La Tabla~\ref{tab:smote_validation} presenta la comparación entre la implementación \textit{dummy-binary} utilizada en esta investigación (que crea un problema binario auxiliar por clase con ruido gaussiano como clase ficticia) y la implementación estándar multiclase de SMOTE (que utiliza máscara binaria objetivo vs resto). Ambas implementaciones invocan internamente el algoritmo \texttt{SMOTE} de la biblioteca \texttt{imbalanced-learn}.

\begin{table}[htbp]
\centering
\caption{Validación de la implementación de SMOTE: por clase (\textit{dummy-binary}) vs estándar multiclase. Macro F1 con regresión logística sobre las 9 combinaciones conjunto $\times$ semilla. Las dos columnas son idénticas porque, fijada la semilla, ambas implementaciones invocan el mismo algoritmo \texttt{imblearn.SMOTE} con los mismos $k$-vecinos intra-clase, generando muestras sintéticas bit-a-bit idénticas.}
\label{tab:smote_validation}
\small
\setlength{\tabcolsep}{4pt}
\begin{tabular}{llrrr}
\toprule
\textbf{Conjunto} & \textbf{Sem.} & \textbf{F1 \textit{dummy}} & \textbf{F1 multic.} & \textbf{$|\Delta|$ (pp)} \\
\midrule
sms\_spam\_10shot              & 42  & 0.8683 & 0.8683 & 0.0000 \\
                               & 123 & 0.8853 & 0.8853 & 0.0000 \\
                               & 456 & 0.8742 & 0.8742 & 0.0000 \\
\midrule
20newsgroups\_10shot           & 42  & 0.7493 & 0.7493 & 0.0000 \\
                               & 123 & 0.7396 & 0.7396 & 0.0000 \\
                               & 456 & 0.7409 & 0.7409 & 0.0000 \\
\midrule
hate\_speech\_davidson\_10shot & 42  & 0.5249 & 0.5249 & 0.0000 \\
                               & 123 & 0.5271 & 0.5271 & 0.0000 \\
                               & 456 & 0.5298 & 0.5298 & 0.0000 \\
\midrule
\multicolumn{2}{l}{\textbf{Total} (27 comparaciones, 3 clasificadores: LR, SVC, Ridge)} & \multicolumn{2}{r}{$|\Delta|$ máximo:} & \textbf{0.0000} \\
\midrule
\multicolumn{5}{l}{\small Umbral de validación: $|\Delta| < 0.1$ pp. \textbf{Resultado: validado.}} \\
\bottomrule
\end{tabular}
\end{table}


La diferencia máxima entre ambas implementaciones es exactamente $0.0000$ puntos porcentuales en las 27 comparaciones realizadas. Este resultado, que a primera vista podría parecer inverosímil, es matemáticamente esperado: con \texttt{random\_state} fijo, el algoritmo SMOTE de \texttt{imbalanced-learn} aplica $k$-vecinos más cercanos exclusivamente sobre las muestras de la clase objetivo. Como en ambas implementaciones el conjunto de embeddings reales de la clase es idéntico, la búsqueda de vecinos produce los mismos pares y los puntos aleatorios de interpolación son idénticos al estar inicializados por la misma semilla. La consecuencia es que las muestras sintéticas generadas son bit-a-bit idénticas, lo que produce conjuntos de entrenamiento aumentados idénticos y, por tanto, clasificadores deterministas con $F1$ idénticos. Cualquier diferencia distinta de cero indicaría un error en la implementación. Este resultado confirma que las ganancias reportadas en esta investigación no se deben a una línea base artificialmente débil.

\section{Análisis de Asimetría de Semillas}
\label{sec:asimetria_semillas}

La Tabla~\ref{tab:seed_asymmetry} presenta el análisis de varianza entre semillas.

\begin{table}[htbp]
\centering
\caption{Impacto de la asimetría de semillas en las pruebas estadísticas (pond. suave vs SMOTE)}
\label{tab:seed_asymmetry}
\small
\begin{tabular}{lrrrrrr}
\toprule
\textbf{Subconjunto} & \textbf{$N$ pares} & \textbf{$\bar{\Delta}$ (pp)} & \textbf{$t$} & \textbf{$p$} & \textbf{$d$ de Cohen} & \textbf{Vic.\%} \\
\midrule
Todos los clasif.       & 105 & +2.25 & 7.63  & $<$0.0001 & 0.745 & 83.8\% \\
Solo estocásticos (MLP+RF)  & 42  & +1.79 & 3.39  & 0.0015    & 0.523 & 78.6\% \\
Determinísticos (LR+SVC+Ridge) & 63  & +2.55 & 7.49  & $<$0.0001 & 0.943 & 87.3\% \\
\midrule
\multicolumn{7}{l}{\small 80\% de los grupos LLM + clasificador determinístico tienen varianza cero entre semillas.} \\
\multicolumn{7}{l}{\small \textbf{El resultado sigue siendo significativo con solo clasificadores estocásticos} ($p=0.0015$).} \\
\bottomrule
\end{tabular}
\end{table}


Las generaciones del LLM se almacenan en caché mediante hash MD5 del prompt. Cuando el mismo prompt se ejecuta con diferentes semillas de muestreo de datos pero el mismo contenido de clase, la caché devuelve respuestas idénticas. Esto produce varianza cero entre semillas para clasificadores determinísticos (regresión logística, SVC, Ridge). El 80\% de las combinaciones LLM + clasificador determinístico presentan esta característica.

Restringiendo el análisis exclusivamente a clasificadores estocásticos (MLP y bosque aleatorio), la ponderación suave mantiene una mejora de +1.79 puntos porcentuales ($p = 0.0015$, $d = 0.52$). El resultado sigue siendo estadísticamente significativo con un tamaño de efecto mediano.

Las pruebas t pareadas siguen siendo válidas porque comparan pares de configuraciones (método vs SMOTE) bajo las mismas condiciones. Los intervalos de confianza podrían subestimar la varianza real al no capturar la variabilidad que introduciría la regeneración de muestras sintéticas con diferentes prompts. Esta limitación se reconoce explícitamente y se discute en el Capítulo~\ref{ch:conclusiones}.
